\subsection{Experimental Setup and Datasets}
\label{sec:setup}

We evaluate SynGS on two widely used novel view synthesis benchmarks: Mip-NeRF 360~\cite{Barron22_MipNeRF360} and LLFF 360~\cite{Mildenhall19_LLFF}. All images are downsampled by a factor of 8. Following sparse-view evaluation protocols, Mip-NeRF 360 scenes are tested with 4, 6, and 9 uniformly sampled input views. The reconstruction quality is evaluated using PSNR, SSIM~\cite{Wang04_SSIM}, and LPIPS~\cite{Zhang18_LPIPS}.

We compare against the baselines reported in Tables~\ref{tab:mip360} and~\ref{tab:llff}: RegNeRF~\cite{Niemeyer22_RegNeRF}, SparseNeRF~\cite{Wang23_SparseNeRF}, 3DGS~\cite{Kerbl23_3DGS}, CoR-GS~\cite{Zhang24_CoRGS}, DropGaussian~\cite{Park25_DropGaussian}, and D$^2$GS~\cite{Song26_D2GS}. All methods are evaluated under consistent input settings within each table.

Our implementation follows the standard 3DGS optimization framework. Depth priors are obtained from VGGT-Depth, i.e., a depth prior adaptation branch trained on a pre-trained VGGT backbone. Consistent with Sec.~\ref{sec:depth_densify}, view- and ray-level selection during densification uses dataset ground-truth depth. The model is trained for 30K iterations, with 3DGS initialization optimized for 10K iterations. Floaters are periodically removed during training. Except for the proposed modules, other training configurations follow the vanilla 3DGS settings. Experiments are conducted using PyTorch on an NVIDIA RTX 4090 GPU.
